% ---------------------------------------------------------------------------
\section{Discussion}

\subsection{What we know, in one page}

The Voynich text is produced by a writing system with three levels of rules:
\begin{enumerate}
\item \textbf{Inside the word and between neighbouring words:} a chain of glyphs with its own sequence rules, in which
  spaces are weak boundaries placed where the chain allows; junction rules between the end of a word and the start of the
  next, applied while writing and looking one word ahead.
\item \textbf{Within the line:} the line is a closed unit with its own starting and ending forms; choices between variants
  follow a short state (two to three words) and drift towards simpler forms from left to right; words are repeated at
  once and the repetition fades within a few words.
\item \textbf{Between lines and across the book:} the line below copies from the line immediately above, adapting the
  copies to the junction rules, and avoids starting like it; the paragraph has its own opening and closing; page,
  bifolium and quire have an identity; the bifolium is a working unit.
\end{enumerate}
Many of these rules belong to the \textbf{system}, not to the person: they hold equally across hands, across Currier's
two ``languages'', in all sections and in all three transliterations.

\subsection{What this excludes and what it leaves open}

\noindent\textbf{It excludes}, with the controls described:
\begin{itemize}
\item a natural language written in an ordinary alphabet, read with a simple or homophonic substitution in one of the
  languages tested;
\item the published readings we could test;
\item the simulated historical encryption mechanisms, taken one at a time;
\item a text invented by hand without rules (the volunteers' gibberish does not have the Voynich's rules);
\item that the Voynich's rules are the ordinary habits of a scribe (79 real scribes do not have them).
\end{itemize}

\noindent\textbf{It leaves open} three scenarios, which our data cannot separate:
\begin{enumerate}
\item \textbf{A language written with a very artificial writing system}: for example a system in which the variants
  (\eva{k/t}, \eva{-ey/-dy}, \ldots) carry no meaning but follow rules of form, and in which spaces do not separate the
  words of the underlying text. A reading would be possible only with an external anchor.
\item \textbf{A code or cipher of an untested kind}, for example with a word repertory or a lost key book. In this case
  statistics is not enough.
\item \textbf{No message:} a text produced by a procedure executable by hand. A procedure of this kind (the
  self-citation generator of Timm and Schinner, with a junction rule added by us) already reproduces many basic properties
  (predictability, similarity between words of the page, repetitions, link between neighbouring words), but not the line
  and page rules described here.
\end{enumerate}
Our data shift the weight slightly towards scenarios 1 and 3 compared with a classical historical cipher, because the
Voynich's rules are ``rules of form'' and no cipher we tested reproduces them. But they \textbf{do not separate} them.

\subsection{A caveat: a message can hide in the choices}

In related work (a tool we call the ``voynichizer'', developed in parallel), we verified that a real text, compressed and
encrypted with a key, can be hidden in the choices between variants (\eva{k/t}, \eva{sh/ch}, \eva{-ey/-dy}, \eva{qo/o},
\ldots) of a text with the properties of the Voynich, \textbf{without making it more recognisable} to a statistical
discriminator. The capacity of this channel is about 0.8 bits per choice, that is about 46{,}000 bits for a book like the
Voynich (the equivalent of about 2{,}000 compressed Latin words). Therefore \textbf{statistical properties alone cannot
prove that the text has no meaning}: a text ``without a message'' and a text with a message hidden in its choices can have
the same statistics.

\subsection{Why statistics alone will not lead to a reading}

All our tools measure \textbf{regularities}. A reading requires knowing what the signs \textbf{represent}, and
regularities do not say this: a very regular writing system can hide a language, a code or nothing. A reading would need
an external element: a text of which the Voynich is a copy (a known herbal or treatise), a secure identification of many
plants to use as cribs, a key found in an archive.

\subsection{What would change our mind}

\begin{itemize}
\item A reading system that passes the controls used here: one that works on the real text and not on the shuffled or
  message-free text, that produces meaningful sentences on pages never seen before, and that explains the writing rules
  described (junction rules, closed line, state of the choices, copying from the line above, margin).
\item A medieval scribe, or a historical cipher, with the same rules: it would change the reading of the whole work.
\item A simple procedure that reproduces all of them together.
\end{itemize}

% ---------------------------------------------------------------------------
\section{Limitations and declared gaps}
\label{sec:limits}

\begin{itemize}
\item \textbf{Small effects.} Many of the new properties are small in absolute terms (agreement of choices is about 13
  percentage points between adjacent words) and required methodological corrections. We report them with intervals and
  replications, but they remain subtle measurements.
\item \textbf{Transliterations.} Everything depends on transliterations: the three used disagree on about one word in
  eight. The main measures hold on all three, but without the images we cannot separate what the scribe did from what the
  transcribers read (for example the \eva{q} dropping after a drawing).
\item \textbf{Comparison scribes.} 79 Nordic and German scribes; Latin or Italian fifteenth-century scribes transcribed
  at facsimile level are missing, as are medieval English scribes (known for very variable spelling). The Nordic and
  German manuscripts do not mark changes of hand; for the state of the choices the 73 German manuscripts were measured
  together, so a window in a single scribe could be diluted.
\item \textbf{Historical ciphers.} Only one real cipher (the Copiale, eighteenth century). The DECODE archive, which
  collects hundreds of fifteenth- and sixteenth-century ciphers, requires an account and has no reuse licence; no
  fifteenth-century cipher appears there with both transcription and key.
\item \textbf{Tests we could not run}, because the measures, tested on synthetic texts, did not distinguish the cases:
  \begin{itemize}
  \item whether the state of the choices lasts a number of \textbf{words} or a span of \textbf{writing time} (glyphs);
  \item whether the states of different choices \textbf{change at the same points} (which would indicate hidden units of a
    few words);
  \item whether the \textbf{short} state crosses the line break (the slow one does).
  \end{itemize}
\item \textbf{A test with few data:} whether the state survives the gap of a drawing (the estimate is positive but the
  interval is too wide to tell).
\item \textbf{Not yet done:} a comparison between manuscripts and printed books of the same corpus.
\end{itemize}

% ---------------------------------------------------------------------------
\section{Reproducibility}

Each experiment has:
\begin{itemize}
\item a dated \textbf{preregistration} (\texttt{preregistrazioni/eNN.md}), written before running it;
\item a \textbf{program} (\texttt{esperimenti/eNN\_*.py}) in Python 3.12, with fixed random seeds;
\item an \textbf{execution record} with the version of code and data (\texttt{risultati/provenienza/});
\item a \textbf{result} as a table (\texttt{risultati/eNN\_*.md}) and in full (\texttt{.json});
\item an entry in the \textbf{laboratory notebook}, written only by adding.
\end{itemize}
Re-running an experiment gives the same numbers. The repository is currently private; copyright-protected data are not
included and must be downloaded again from the sources (Appendix~D).
